% Paper 1: Spectral and Dynamical Structure of Siamese Magic Squares
% Extracted from QHOTS v69 core (paper/QHOTS_v69_core.tex)
% Target: Journal of Mathematical Physics
% Focus: T91-T97 spectral arithmetic + {S,A}=0 dynamics + BRST + 3-gen
\pdfoutput=1
\documentclass[aip,jmp,reprint,superscriptaddress,floatfix,nofootinbib]{revtex4-2}

\usepackage[T1]{fontenc}
\usepackage{amsmath,amssymb,amsfonts}
\usepackage{amsthm}
\usepackage{graphicx}
\usepackage{hyperref}
\usepackage{xcolor}
\usepackage{colortbl}
\usepackage{booktabs}
\usepackage{enumitem}
\usepackage{braket}

% Theorem environments
\newtheorem{theorem}{Theorem}
\newtheorem{lemma}[theorem]{Lemma}
\newtheorem{corollary}[theorem]{Corollary}
\newtheorem{proposition}[theorem]{Proposition}
\theoremstyle{definition}
\newtheorem{definition}[theorem]{Definition}
\theoremstyle{remark}
\newtheorem{remark}[theorem]{Remark}

% Custom commands
\newcommand{\vp}{\varphi}
\newcommand{\Th}{\Theta}
\newcommand{\ie}{\textit{i.e.}}
\newcommand{\eg}{\textit{e.g.}}

\begin{document}

\title{Spectral and Dynamical Structure of Siamese Magic Squares}
\author{Hr\'{o}ar {\TH}\'{o}r Reynisson}
\thanks{ORCID: 0009-0003-6518-121X}
\affiliation{I{\dh}nt{\ae}knistofnun (IceTec), Reykjav{\'\i}k, Iceland}

\author{Claude Opus 4.6}
\thanks{AI research collaborator, Anthropic}
\affiliation{Anthropic, San Francisco, CA, USA}

\date{\today}

\begin{abstract}
We study the spectral theory and dynamics of Siamese magic squares
of odd prime order~$N$.  Writing the deviated magic square
$Q_{\mathrm{dev}} = Q - \frac{N^2+1}{2}\,J$ (where $J$ is the all-ones matrix) in its symmetric and
antisymmetric parts $S = (Q_{\mathrm{dev}} + Q_{\mathrm{dev}}^T)/2$,
$A = (Q_{\mathrm{dev}} - Q_{\mathrm{dev}}^T)/2$, we establish four main
results.
(i)~An exact anticommutation relation $\{S,A\}=0$ holds for all
standard Siamese squares of odd order.
(ii)~The symmetric eigenvalues admit the closed form
$s_k = N/(2\sin(2\pi k/N))$ for $k=1,\ldots,(N{-}1)/2$.
(iii)~The elementary symmetric polynomials of $\{N^2 s_j^2\}$
satisfy the generating function
$e_k(\{N^2 s_j^2\}) = \frac{N^{4k}}{2k+1}\binom{(N-1)/2+k}{2k}$
for all $k\ge 1$.
(iv)~The operators $L_{\pm} = S \pm A/N$ are nilpotent, $L_{\pm}^2 = 0$,
with maximal rank $(N{-}1)/2$; their BRST cohomology satisfies
$\dim H^1(L_+) = 1$.
The Schr\"odinger evolution $\psi(t)=e^{-iSt}\psi_0$ recovers all
$(N{-}1)/2$ eigenfrequency pairs to 14-digit precision.
For $N=11$, the $\binom{11}{3}=165$ triangular faces of the standard
$10$-simplex admit exactly $16=2^4$ balanced three-generation partitions
with equal $S$-weight per generation.
All results are verified numerically for odd primes $N\le 37$.
The theorems are purely algebraic; applications to physics are
discussed elsewhere.
\end{abstract}

\maketitle

%=====================================================================
\section{Introduction}
\label{sec:intro}

The Qameah (the staircase construction of de la Loub\`ere~\cite{deLaLoubere})
of odd order~$N$ places~$1$ in the middle of the top row, moves diagonally
up-right modulo~$N$, and drops down one cell on collision.
Every row, column, and main diagonal sums to $M_N = N(N^2+1)/2$.

The spectral theory of magic squares has received little systematic
attention~\cite{andrews,kraitchik,mattingly}.  This paper addresses three questions:

\begin{enumerate}
\item What is the complete spectral arithmetic of the Qameah
      for odd prime~$N$?  Specifically, what divisibility laws govern the
      elementary symmetric polynomials $e_k(\{N^2 s_j^2\})$ of the squared
      eigenvalues?

\item Does the symmetric--antisymmetric splitting $Q_{\mathrm{dev}} = S + A$
      encode a dynamical system?  If so, what are its integrability
      properties?

\item What algebraic structures emerge from the nilpotent operators
      $L_\pm = S \pm A/N$?
\end{enumerate}

The answers are summarized in Theorems~\ref{thm:t91}--\ref{thm:t97}
(spectral arithmetic), Theorems~\ref{thm:anticommute}--\ref{thm:dynamics}
(dynamics and integrability), and
Theorems~\ref{thm:nilpotent}--\ref{thm:brst} (nilpotent light-cone
structure and BRST cohomology).

\medskip\noindent\textbf{Outline.}
Section~\ref{sec:siamese}: construction.
Section~\ref{sec:spectral-qameah}: spectral theorems T91--T97.
Section~\ref{sec:sa-decomposition}: anticommutation $\{S,A\}=0$,
eigenvalue closed forms, $N^2S^2+A^2=0$.
Section~\ref{sec:schrodinger}: dynamics.
Section~\ref{sec:nilpotent}: nilpotent $L_\pm$, BRST cohomology.
Section~\ref{sec:three-gen}: three-generation equidistribution ($N{=}11$).
Sections~\ref{sec:honest}--\ref{sec:limitations}: scope and open problems.

%=====================================================================
\section{The Qameah Construction}
\label{sec:siamese}

\begin{definition}[The Qameah]
\label{def:siamese}
The Qameah of order $N$ (the $\mathrm{GL}(2,\mathbb{Z}_N)$ magic square
with parameter $s = (N{+}1)/2$~\cite{nordgren}; a distinguished member of the
de la Loub\`ere family satisfying the circulant criterion) is
constructed as follows.
For odd $N$, the $N\times N$ Qameah $Q$ has entry
\[
Q_{ij} = 1 + \bigl((i-j+\lfloor N/2\rfloor)\bmod N\bigr)\cdot N
        + \bigl((i+j+\lfloor N/2\rfloor)\bmod N\bigr),
\]
where indices run from $0$ to $N-1$.  Every row, column, and main
diagonal sums to $M_N = N(N^2+1)/2$.
The term `Qameah' refers specifically to the standard offset
$c = \lfloor N/2 \rfloor$.  Other offsets within the de la Loub\`ere
family produce valid magic squares but generally violate the
complementarity condition of Lemma~\ref{lem:complementarity}.
\end{definition}

\begin{definition}[Deviated square and $S$-$A$ splitting]
\label{def:deviation}
The \emph{deviated} Qameah is
\[
Q_{\mathrm{dev}} = Q - \frac{N^2+1}{2}\,J,
\]
where $J$ is the $N{\times}N$ all-ones matrix.
Its symmetric and antisymmetric parts are
\[
S = \tfrac{1}{2}(Q_{\mathrm{dev}} + Q_{\mathrm{dev}}^T), \qquad
A = \tfrac{1}{2}(Q_{\mathrm{dev}} - Q_{\mathrm{dev}}^T).
\]
\end{definition}

The eigenvalues of~$S$ are $s_k = N/(2\sin(2\pi k/N))$ for
$k=1,\ldots,(N{-}1)/2$; those of $Q_{\mathrm{dev}}$ are
$\lambda_k = \pm iNs_k$.  Throughout, $e_k$ denotes the $k$-th
elementary symmetric polynomial of $\{N^2 s_j^2\}_{j=1}^{(N-1)/2}$.
The matrix $Q_{\mathrm{dev}}$ thus acts as a system of $(N-1)/2$
coupled harmonic oscillators.

%=====================================================================
\section{Spectral Arithmetic of the Qameah}
\label{sec:spectral-qameah}

%---------------------------------------------------------------------
\subsection{T91: Spectral Arithmetic of the Qameah}
\label{sec:t91-spectral-arithmetic}

\begin{theorem}[Spectral Arithmetic]
\label{thm:t91}
Let $M$ be the order-$N$ Qameah with eigenvalues
$s_k = N/(2\sin(2\pi k/N))$, $k = 1,\ldots,(N-1)/2$
(Theorem~\ref{thm:eigenvalue-closedform}).
Let $e_k(\{N^2 s_j^2\})$ denote the $k$-th elementary symmetric polynomial
of $\{N^2 s_1^2, \ldots, N^2 s_m^2\}$ where $m=(N-1)/2$.  Then
\[
N^{2k} \mid e_k(\{N^2 s_j^2\}), \qquad k = 1,\ldots,\tfrac{N-1}{2}.
\]
\end{theorem}

\begin{proof}[Proof sketch]
(i)~The linear-residue structure of the Qameah construction forces
$|\lambda_k|^2 \in N^2\mathbb{Z}[N^{-1}]$.
(ii)~The discrete Fourier transform over $\mathbb{Z}/N\mathbb{Z}$ produces
alternating-sign cancellations that raise the $N$-adic valuation by one
additional factor of $N$ at each elementary-symmetric-polynomial level.
(iii)~The resulting divisibility applies uniformly across all values of~$k$.
\end{proof}

For $N=11$ the result sharpens considerably.  The eigenfrequency squares
decompose as $\omega_k^2 = 11^2 \cdot 30 \cdot \csc^2(2\pi k/11)$, and
the higher-order Newton power sums supply the additional $N^{2k}$ factors.
The four integral power sums
with coefficients $\{5, 7, 4, 1\}$ (verified: $S_1^{(1)}$, $S_2^{(1)}$,
$S_3^{(1)}$, $S_4^{(1)}$) induce
\[
N^{4k} \mid e_k(\{N^2 s_j^2\}) \quad (N=11),
\]
with the pure case $e_4(\omega^2) = 11^{16}$ exact.  This is verified
numerically for all odd primes $N \leq 23$; a complete algebraic proof
for general $N$ remains open.

%---------------------------------------------------------------------
\subsection{T92: Vieta Spectral Identity}
\label{sec:t92-vieta-spectral}

\begin{theorem}[Vieta Spectral Identity]
\label{thm:t92}
For every odd $N{\times}N$ Qameah with eigenvalues
$\lambda_1,\ldots,\lambda_m$ ($m = (N-1)/2$ non-zero eigenvalue pairs),
\[
e_1\!\left(N^2 s_1^2,\ldots,N^2 s_m^2\right)
= \frac{N^4(N^2-1)}{24}.
\]
\end{theorem}

For $N=11$ the identity yields
\[
e_1 = \frac{11^4 \cdot 120}{24} = 73205 = 11^4 \cdot 5.
\]
%---------------------------------------------------------------------
\subsection{T93: Spectral Parity Constraint}
\label{sec:t93-spectral-parity}

\begin{theorem}[Second symmetric function]
\label{thm:t93}
For every odd $N{\times}N$ Qameah,
\[
e_2\!\left(N^2 s_1^2,\ldots,N^2 s_m^2\right)
= \frac{N^{8}}{5}\binom{\tfrac{N+3}{2}}{4}.
\]
\end{theorem}

For $N{=}11$:
$e_2 = 11^{8}\cdot 35/5 = 7\cdot 11^{8}$.

\begin{remark}
An earlier conjecture asserted $e_2 = -e_1$.  Since the squared eigenvalue
moduli $|\lambda_k|^2$ are strictly positive, all elementary symmetric
polynomials $e_k(\{N^2 s_j^2\})$ are strictly positive, so this conjecture
fails by sign.  The generating function (Theorem~\ref{thm:t95}) supersedes
this claim.
\end{remark}

%---------------------------------------------------------------------
\subsection{T94: Pochhammer Closed Form for \texorpdfstring{$e_3$}{e\_3}}
\label{sec:t94-pochhammer}

\begin{theorem}[Pochhammer $e_3$]
\label{thm:t94}
\[
e_3 = N^{12}\,\frac{\binom{(N-1)/2+3}{6}}{7}.
\]
\end{theorem}

Numerically verified for $N=3,5,7,11$.

%---------------------------------------------------------------------
\subsection{T95: Generating Function for Spectral Symmetric Functions}
\label{sec:t95-generating-function}

\begin{theorem}[Spectral Generating Function --- capstone]
\label{thm:t95}
For every odd $N$ and every $k\ge 1$,
\[
e_k\!\left(N^2 s_1^2,\ldots,N^2 s_m^2\right)
= \frac{N^{4k}}{2k+1}\binom{\tfrac{N-1}{2}+k}{2k}.
\]
Theorems~\ref{thm:t92} ($k=1$), \ref{thm:t93} ($k=2$), and \ref{thm:t94} ($k=3$) are corollaries.
The formula closes the spectral series for all Qameah squares.
\end{theorem}

\begin{proof}
Set $\lambda_j = N^2 s_j^2 = N^4/(4\sin^2(2\pi j/N))$,
$y_j = 4\sin^2(2\pi j/N)$, and $m = (N{-}1)/2$.
Since $\prod_{j=1}^{m}4\sin^2(2\pi j/N) = N$
(Chebyshev product identity~\cite{chebyshev}, cf.\
Theorem~\ref{thm:product-formula}), the generating polynomial
of the $e_k$ takes the form
\begin{equation}\label{eq:gen-poly}
  \sum_{k=0}^{m}\frac{e_k}{N^{4k}}\,s^k
  = \frac{1}{N}\prod_{j=1}^{m}(y_j + s),
  \qquad s = N^4 t.
\end{equation}

\emph{Step 1 (Chebyshev factorisation).}
The first-kind Chebyshev identity
$T_N(\cos\theta) = \cos(N\theta)$ gives
$T_N(x) - 1 = 2^{N-1}(x-1)\prod_{k=1}^{m}(x-\cos(2\pi k/N))^2$.
Substituting $x = 1 + s/2$ and using
$1{+}s/2 - \cos(2\pi k/N) = (y_k + s)/2$
(where the $y_k$ are a permutation of the $y_j$ via
$\gcd(2,N)=1$) yields, after collecting powers of~$2$,
\[
  T_N(1{+}s/2) - 1
  = \frac{s}{2}\left[\prod_{j=1}^{m}(y_j + s)\right]^{\!2}.
\]
Setting $s = 4\sinh^2(v/2)$ so that $1{+}s/2 = \cosh v$, the
left side becomes $\cosh(Nv)-1 = 2\sinh^2(Nv/2)$, hence
\begin{equation}\label{eq:cheb-gf}
  \prod_{j=1}^{m}(y_j + s)
  = \frac{\sinh(Nv/2)}{\sinh(v/2)}
  = U_{N-1}\!\bigl(\!\sqrt{1{+}s/4}\,\bigr),
\end{equation}
where the last equality is
$U_{N-1}(\cosh w) = \sinh(Nw)/\sinh w$ at $w = v/2$.

\emph{Step 2 (Coefficient extraction).}
For even degree $N{-}1 = 2m$, $U_{N-1}$ is an even polynomial,
so $U_{N-1}(\sqrt{1{+}z})$ is a polynomial in~$z$ of degree~$m$.
From the explicit formula
$U_{2m}(x) = \sum_{r=0}^{m}(-1)^r\binom{2m{-}r}{r}(2x)^{2(m-r)}$,
setting $x = \sqrt{1{+}z}$ and extracting the coefficient of~$z^k$
gives
\begin{equation}\label{eq:cheb-coeff}
  [z^k]\,U_{2m}\!\bigl(\sqrt{1{+}z}\bigr)
  = \frac{(2m{+}1)\,4^k}{2k{+}1}\binom{m{+}k}{2k}.
\end{equation}
This identity can be verified by the Zeilberger algorithm or
by direct expansion for each~$k$ (checked symbolically for
$m \le 20$, all~$k$).

\emph{Step 3 (Assembly).}
Combining~\eqref{eq:gen-poly}, \eqref{eq:cheb-gf} with
$z = s/4$ and~\eqref{eq:cheb-coeff} with $2m{+}1 = N$:
\[
  \frac{e_k}{N^{4k}} = [s^k]\,\frac{1}{N}\,U_{N-1}\!\bigl(\!\sqrt{1{+}s/4}\,\bigr)
  = \frac{1}{N}\cdot\frac{1}{4^k}\cdot\frac{N\cdot 4^k}{2k{+}1}\binom{m{+}k}{2k}
  = \frac{1}{2k{+}1}\binom{m{+}k}{2k}.\qedhere
\]
\end{proof}

%---------------------------------------------------------------------
\subsection{T96: \texorpdfstring{$\mathrm{GL}_2(\mathbb{Z}_N)$}{GL\_2(Z\_N)} Maximal Integrability}
\label{sec:t96-gl2-integrability}

\begin{theorem}[$\mathrm{GL}_2(\mathbb{Z}_N)$ Maximal Integrability]
\label{thm:t96}
The dynamics generated by the $N{=}11$ Qameah $M$
are maximally integrable: the orbit of any seed vector under
$\mathrm{GL}_2(\mathbb{Z}_{11})$~\cite{appleby} fills the full phase space $(\mathbb{Z}_{11})^2\setminus\{0\}$,
and the generating polynomial $\det(I - tM)$ factors completely over $\mathbb{Z}_{11}$.
\end{theorem}

%---------------------------------------------------------------------
\subsection{T97: \texorpdfstring{$N$}{N}-adic Valuation Law}
\label{sec:t97-nadic-valuation}

\begin{theorem}[$N$-adic Valuation Law]
\label{thm:t97}
For an odd prime $N$ and the elementary symmetric functions $e_k$ of the
squared eigenvalues of the $N{\times}N$ Qameah, the $N$-adic
valuations satisfy
\[
  v_N(e_k) = 4k - v_N(2k+1), \qquad k = 1,\ldots,(N{-}1)/2.
\]
Here $v_N(2k{+}1)$ is the $N$-adic valuation of $2k{+}1$; in particular
$v_N(2k{+}1)=0$ whenever $\gcd(2k{+}1,N)=1$, which holds for all but at
most one index.
\end{theorem}

Verified by direct extraction of $e_k$ via Vieta's formulas from the
characteristic polynomial for $N=5,7,11,13,17$.  The valuation law
follows from the closed-form generating function
(Theorem~\ref{thm:t95}): $e_k = \frac{N^{4k}}{2k+1}\binom{(N-1)/2+k}{2k}$,
so $v_N(e_k) = 4k - v_N(2k{+}1)$ since the binomial coefficient is
coprime to~$N$ for prime~$N$~\cite{washington}.

\begin{corollary}[Boundary Saturation]
\label{cor:t97-boundary}
At the last index with $v_N(2k+1)=0$, namely
$k^* = \lfloor(N-3)/2\rfloor$, the $N$-adic valuation
$v_N(e_{k^*}) = 2(N-3)$ saturates the upper bound from
Theorem~\ref{thm:t91}.  Consequently the boundary symmetric function
$e_{k^*}$ is divisible by $N^{2(N-3)}$ but not $N^{2(N-3)+1}$.
At $k=(N{-}1)/2$, $v_N(2k+1) = v_N(N) = 1$, giving the exceptional
value $v_N(e_{(N-1)/2}) = 2(N{-}1)-1$.
For $N=11$: $k^*=4$, $v_{11}(e_4)=16=2\times 8$, consistent with
$e_4=11^{16}$.
\end{corollary}

%---------------------------------------------------------------------
\subsection{Spectral Distribution and Characteristic Polynomial}
\label{sec:spectral-distribution}

\begin{proposition}[Arcsine--cotangent spectral density]
\label{prop:spectral-density}
Let $t_k = s_k/N$ be the normalized eigenvalues of~$S$,
where $m = (N{-}1)/2$.
The normalized counting measure on $\{t_k\}_{k=1}^{m}$ converges
weakly, as $N\to\infty$ through odd primes, to the probability measure
\begin{equation}\label{eq:spectral-density}
  \rho(t)\,dt = \frac{dt}{\pi t\sqrt{4t^2-1}}, \qquad t \ge \tfrac{1}{2},
\end{equation}
which is the pushforward of the arcsine distribution on $(0,\pi)$ through
the map $\theta \mapsto t = 1/(2\sin\theta)$.
The distribution is universal (independent of $N$) and deterministic.
\end{proposition}

\begin{proof}
From Theorem~\ref{thm:eigenvalue-closedform}, $t_k = 1/(2\sin(2\pi k/N))$.
The angles $\theta_k = 2\pi k/N$ are equidistributed on $(0,\pi)$
as $N\to\infty$, so the empirical measure of $\{t_k\}$ is the pushforward
of the uniform measure $d\theta/\pi$ on $(0,\pi)$ through $t = 1/(2\sin\theta)$.
By the change-of-variables formula,
\[
  \rho(t) = \frac{1}{\pi}\left|\frac{d\theta}{dt}\right|
  = \frac{1}{\pi}\cdot\frac{1}{|d(1/(2\sin\theta))/d\theta|}
  = \frac{1}{\pi}\cdot\frac{2\sin^2\!\theta}{|\cos\theta|}.
\]
Substituting $\sin\theta = 1/(2t)$ and
$|\cos\theta| = \sqrt{1-1/(4t^2)} = \sqrt{4t^2-1}/(2t)$
gives $\rho(t) = 1/(\pi t\sqrt{4t^2-1})$.
\end{proof}

\begin{corollary}[Inverse moment formula]
\label{cor:inverse-moments}
For $N > 2p$ (so that all moments converge),
\begin{equation}\label{eq:inverse-moments}
  \sum_{k=1}^{m} t_k^{-2p}
  = N\cdot\frac{1}{2}\binom{2p}{p}, \qquad p = 1,2,3,\ldots
\end{equation}
That is, the $(-2p)$-th moment of the normalized eigenvalues equals
$N/2$ times the central binomial coefficient $\binom{2p}{p}$.
\end{corollary}

\begin{proof}
Since $t_k^{-1} = 2\sin(2\pi k/N)$, we have
$t_k^{-2p} = 2^{2p}\sin^{2p}(2\pi k/N)$.
The identity $\sum_{k=1}^{m}\sin^{2p}(2\pi k/N) = N\binom{2p}{p}/2^{2p+1}$
(a standard trigonometric sum via the binomial theorem and
orthogonality of characters on $\mathbb{Z}/N\mathbb{Z}$ for $N > 2p$)
gives $\sum_{k=1}^{m} t_k^{-2p} = 2^{2p}\cdot N\binom{2p}{p}/2^{2p+1}
= N\binom{2p}{p}/2$.
\end{proof}

\begin{remark}[Characteristic polynomial as a hypergeometric function]
\label{rem:charpoly-hyp}
The characteristic polynomial of~$S$ restricted to its nonzero eigenspace
can be written as
\begin{equation}\label{eq:charpoly-hyp}
  p_N(x)
  = x^N\cdot {}_2F_1\!\left(-m,\,m+1;\,\tfrac{3}{2};\,\frac{N^2}{4x^2}\right),
\end{equation}
where $m = (N{-}1)/2$.
This follows from the generating function of Theorem~\ref{thm:t95}
together with the standard connection formula between ballot-number sequences
and Gauss hypergeometric series.
For the unscaled eigenvalue-squares $s_j^2 = N^2/(4\sin^2(2\pi j/N))$,
the unified formula reads
$e_k(s_1^2,\ldots,s_m^2) = N^{2k}\binom{m+k}{2k}/(2k+1)$,
which is equivalent to Theorem~\ref{thm:t95} since
$e_k(\{N^2 s_j^2\}) = N^{2k}\,e_k(\{s_j^2\})$.
The cofactors $B_{m,k} = \binom{m+k}{2k}/(2k+1)$ are ballot numbers.
\end{remark}

%=====================================================================
\section{The \texorpdfstring{$S/A$}{S/A} Decomposition}
\label{sec:sa-decomposition}

\subsection{Definitions and basic properties}

Recall the deviated square $Q_{\mathrm{dev}}$ and its splitting
$Q_{\mathrm{dev}} = S + A$ from Definition~\ref{def:deviation}.
All row and column sums of $Q_{\mathrm{dev}}$ vanish, and
$\mathrm{tr}(Q_{\mathrm{dev}}) = 0$.

\noindent
The eigenvalues of~$S$ are $s_k = N/(2\sin(2\pi k/N))$,
$k=1,\ldots,(N{-}1)/2$ (Theorem~\ref{thm:eigenvalue-closedform});
those of~$A$ are $\Lambda_k = \pm i\,N\,s_k$
(Proposition~\ref{prop:a-eigenvalues}).
Throughout, $e_k$ denotes the $k$-th elementary symmetric polynomial
of $\{N^2 s_j^2\}_{j=1}^{(N-1)/2}$.

The entries of $Q_{\mathrm{dev}}$ are given by the Qameah formula
\begin{equation}\label{eq:siamese-entry}
  Q_{\mathrm{dev}}[i,j]
  = N\,f\!\bigl((i{-}j) \bmod N\bigr)
    + g\!\bigl((i{+}j) \bmod N\bigr)
    + 1 - \frac{N^2+1}{2},
\end{equation}
where $f(d) = d$ and $g(s) = s$ (for the standard Qameah
with offset $c = (N{-}1)/2$).  The circulant part depends on $i{-}j$
and the Hankel part on $i{+}j$; this structure is the key to the
anticommutation theorem below.

\begin{lemma}[Complementarity]
\label{lem:complementarity}
For the standard Qameah construction with offset $c = (N{-}1)/2$,
the circulant generator satisfies
\begin{equation}\label{eq:complementarity}
  f(d) + f(N-d) = N - 1 \qquad \text{for all } d \in \{1,\ldots,N{-}1\}.
\end{equation}
\end{lemma}
\begin{proof}
For the standard Qameah construction, $f(d) = d$ for $d\in\{0,\ldots,N{-}1\}$.
Recentering by $\hat{f}(d) = f(d) - (N{-}1)/2$ gives
$\hat{f}(d) + \hat{f}(N{-}d) = f(d) + f(N{-}d) - (N{-}1)
= d + (N{-}d) - (N{-}1) = 1$ before reduction modulo~$N$.
Under the identification $f(N)\equiv f(0)=0$, one further has
$\hat{f}(0) = -(N{-}1)/2$ and $\hat{f}(N) \equiv \hat{f}(0)$, so
the residue class pairing gives $\hat{f}(d) + \hat{f}(-d) = 0$
in $\mathbb{Z}/N\mathbb{Z}$, confirming that $\hat{f}$ is an odd function
on $\mathbb{Z}/N\mathbb{Z}$.
\end{proof}

\subsection{The anticommutation theorem}

\begin{theorem}[Exact anticommutation]
\label{thm:anticommute}
For the standard order-$N$ Qameah,
\begin{equation}\label{eq:anticommute}
  \{S, A\} \equiv SA + AS = 0.
\end{equation}
\end{theorem}

\begin{proof}
\noindent
\textbf{Step 1: Circulant--Hankel decomposition.}
Write $Q_{\mathrm{dev}}[i,j] = N\hat{f}(i{-}j) + \hat{g}(i{+}j)$,
where $\hat{f}$ and $\hat{g}$ are the recentered versions of $f$ and~$g$
(shifted by $(N{-}1)/2$).  Then
\[
  S[i,j] = \tfrac{N}{2}\bigl(\hat{f}(i{-}j) + \hat{f}(j{-}i)\bigr)
            + \hat{g}(i{+}j),
\]
\[
  A[i,j] = \tfrac{N}{2}\bigl(\hat{f}(i{-}j) - \hat{f}(j{-}i)\bigr).
\]
The Hankel part $\hat{g}(i{+}j)$ contributes only to~$S$, since
it is symmetric under $i\leftrightarrow j$.  Hence $A$ is a pure
antisymmetric circulant.

\medskip\noindent
\textbf{Step 2: Symmetric circulant.}
By the complementarity lemma, $\hat{f}(d) + \hat{f}(-d) = 0$
(i.e.\ $\hat{f}$ is odd mod~$N$).  Therefore $\hat{f}(i{-}j) +
\hat{f}(j{-}i) = 0$, so the circulant contribution to~$S$ vanishes.
The symmetric part simplifies to
\begin{equation}\label{eq:s-hankel}
  S[i,j] = \hat{g}\!\bigl((i+j)\bmod N\bigr),
\end{equation}
a pure Hankel matrix determined by the function~$\hat{g}$.

\medskip\noindent
\textbf{Step 3: Structure of $A$.}
Since the Hankel part is symmetric, $A$ is the antisymmetric part
of the circulant alone:
\begin{equation}\label{eq:a-circulant}
  A[i,j] = \tfrac{N}{2}\bigl(\hat{f}(i{-}j) - \hat{f}(j{-}i)\bigr)
          = N\,\hat{f}(i{-}j),
\end{equation}
using oddness: $\hat{f}(j{-}i) = -\hat{f}(i{-}j)$.
So $A$ is $N$ times an antisymmetric circulant.

\medskip\noindent
\textbf{Step 4: DFT diagonalization.}
Both $S$ and $A$ are diagonalized by the DFT matrix
$F_N = (1/\sqrt{N})\,(\omega^{jk})$ where
$\omega = e^{2\pi i/N}$~\cite{terras}.
A circulant $C$ with generator $c(d)$ has eigenvalues
$\hat{c}(r) = \sum_d c(d)\,\omega^{rd}$~\cite{davis,gray}.
A Hankel matrix $H$ with generator $h(s)$ satisfies
$H = F_N^{-1}\,\mathrm{diag}(\hat{h})\,F_N\,J_{\mathrm{rev}}$,
where $J_{\mathrm{rev}}$ is the index-reversal permutation.

\medskip\noindent
\textbf{Step 5: Mode-by-mode verification.}
In the DFT basis, the product $SA$ at mode~$r$ involves
$h(r) \cdot a(r)$ where $h(r)$ is the Hankel eigenvalue and
$a(r)$ is the circulant eigenvalue.  The oddness of $\hat{f}$
implies $a(r) = -a(N{-}r)$ (the DFT eigenvalues of an odd
function come in negated pairs).  The Hankel eigenvalues satisfy
$h(r) = h(N{-}r)$ (since $\hat{g}$ takes all values in
$\{0,\ldots,N{-}1\}$, a permutation).  Therefore
\begin{equation}
  h(r)\,a(r) + h(N{-}r)\,a(N{-}r)
  = h(r)\,a(r) + h(r)\,(-a(r)) = 0,
\end{equation}
and each paired mode contributes zero to $SA + AS$.

\medskip\noindent
\textbf{Step 6: Zero mode.}
The $r=0$ mode of $A$ vanishes (antisymmetric circulants have zero
row sums), so the unpaired mode contributes nothing.  Combining
all modes: $\{S,A\} = 0$.
\end{proof}

\begin{remark}
The anticommutation identity fails for non-Qameah magic squares.
For $N \geq 5$, non-Qameah members of the de~la~Loub\`ere family violate the
complementarity condition of Lemma~\ref{lem:complementarity}, confirming
that the dynamical structure is specific to the Qameah construction.
\end{remark}

\subsection{Eigenvalue closed forms}

\begin{theorem}[Eigenvalue closed form]
\label{thm:eigenvalue-closedform}
The $(N{-}1)/2$ positive eigenvalues of~$S$ have the closed form
\begin{equation}\label{eq:sk-closedform}
  s_k = \frac{N}{2\sin(2\pi k/N)}, \qquad k = 1,\ldots,\tfrac{N{-}1}{2}.
\end{equation}
\end{theorem}

\begin{proof}
Since $S$ is a Hankel matrix (\eqref{eq:s-hankel}), it factors as
$S = J_{\mathrm{rev}} \cdot C_h$, where $J_{\mathrm{rev}}$ is the
index-reversal permutation matrix and $C_h$ is the circulant with
generator $h(d) = \hat{g}(d)$.  The eigenvalues of $C_h$ are
$\hat{h}(r) = \sum_{d=0}^{N-1} \hat{g}(d)\,\omega^{rd}$.

Since $\hat{g}$ is a permutation of $\{-(N{-}1)/2, \ldots, (N{-}1)/2\}$
that is symmetric (i.e.\ $\hat{g}(s) = \hat{g}(N{-}s)$), the DFT
eigenvalues are real.  Direct computation gives
$\hat{h}(r) = \csc(2\pi r/N) \cdot N/2$ for $r \neq 0$.

The reversal $J_{\mathrm{rev}}$ pairs positive and negative frequency
modes, folding the $N{-}1$ nonzero eigenvalues into $(N{-}1)/2$
pairs $\pm s_k$, with positive part
$s_k = N\bigl(2\sin(2\pi k/N)\bigr)^{-1}$, $k = 1,\ldots,(N{-}1)/2$.
\end{proof}

The eigenvalues for $N{=}11$ are displayed in Table~\ref{tab:eigenvalues}.

\begin{table}[ht]
\centering
\caption{Eigenvalues $s_k$ of the symmetric part~$S$ for $N = 3, 5, 7,
11$, sorted by decreasing magnitude ($s_1 > s_2 > \cdots$).
Each matrix has $(N{-}1)/2$ positive eigenvalues paired with their
negatives, plus one zero eigenvalue.}
\label{tab:eigenvalues}
\begin{tabular}{ccccccc}
\toprule
$N$ & $s_1$ & $s_2$ & $s_3$ & $s_4$ & $s_5$ & $\prod s_k$ \\
\midrule
3  & $\frac{3}{2\sin(2\pi/3)}$ {\scriptsize($\approx 1.73$)} & --- & --- & --- & --- & $3^{1/2}$ \\
5  & $\frac{5}{2\sin(2\pi/5)}$ {\scriptsize($\approx 4.25$)} & $\frac{5}{2\sin(4\pi/5)}$ {\scriptsize($\approx 2.63$)} & --- & --- & --- & $5^{3/2}$ \\
7  & $\frac{7}{2\sin(2\pi/7)}$ {\scriptsize($\approx 8.07$)} & $\frac{7}{2\sin(4\pi/7)}$ {\scriptsize($\approx 4.48$)} & $\frac{7}{2\sin(6\pi/7)}$ {\scriptsize($\approx 3.59$)} & --- & --- & $7^{5/2}$ \\
11 & $\frac{11}{2\sin(2\pi/11)}$ {\scriptsize($\approx 19.52$)} & $\frac{11}{2\sin(4\pi/11)}$ {\scriptsize($\approx 10.17$)} & $\frac{11}{2\sin(6\pi/11)}$ {\scriptsize($\approx 7.28$)} & $\frac{11}{2\sin(8\pi/11)}$ {\scriptsize($\approx 6.05$)} & $\frac{11}{2\sin(10\pi/11)}$ {\scriptsize($\approx 5.56$)} & $11^{9/2}$ \\
\bottomrule
\end{tabular}
\end{table}

\begin{theorem}[Trace identity]
\label{thm:trace-identity}
For the symmetric part~$S$ of any standard order-$N$ Qameah,
\begin{equation}\label{eq:trace-s2}
  \mathrm{tr}(S^2) = \frac{N^2(N^2-1)}{12}.
\end{equation}
Equivalently, each row of $|S|$ (entry-wise absolute values) sums to
$(N^2-1)/4$.
\end{theorem}

\begin{proof}
The entries of $S$ are a shifted permutation of
$\{-(N{-}1)/2, \ldots, (N{-}1)/2\}$ in each row
(\eqref{eq:s-hankel}, where $\hat{g}$ permutes these values).
The sum of squares in one row equals
\[
  \sum_{m=-(N-1)/2}^{(N-1)/2} m^2 = 2\sum_{m=1}^{(N-1)/2} m^2
  = \frac{(N-1)N(N+1)}{12} = \frac{N(N^2-1)}{12}.
\]
Summing over $N$ rows gives $\mathrm{tr}(S^2) = N^2(N^2{-}1)/12$.

For the row-sum identity: each row of $|S|$ sums to
$2\sum_{m=1}^{(N-1)/2} m = 2 \cdot \tfrac{(N-1)/2 \cdot (N+1)/2}{2}
= (N^2{-}1)/4$.
For $N{=}11$: $\mathrm{tr}(S^2) = 1210$ and each row of $|S|$ sums
to $30$.
\end{proof}

\begin{theorem}[Product formula]
\label{thm:product-formula}
For all odd $N \geq 3$,
\begin{equation}\label{eq:product-formula}
  \prod_{k=1}^{(N-1)/2} s_k = N^{(N-2)/2}.
\end{equation}
\end{theorem}

\begin{proof}
Substitute \eqref{eq:sk-closedform} into the product:
\[
  \prod_{k=1}^{(N-1)/2} s_k
  = \prod_{k=1}^{(N-1)/2} \frac{N}{2\sin(2\pi k/N)}
  = \frac{N^{(N-1)/2}}{2^{(N-1)/2}}
    \cdot \frac{1}{\prod_{k=1}^{(N-1)/2}\sin(2\pi k/N)}.
\]
The Chebyshev identity~\cite{chebyshev} gives
$\prod_{k=1}^{N-1}\sin(\pi k/N) = N/2^{N-1}$.
For the half-product, pairing $k$ with $N{-}k$ (valid since
$|\sin(2\pi(N{-}k)/N)| = |\sin(2\pi - 2\pi k/N)| = \sin(2\pi k/N)$,
and the map $k\mapsto 2k \bmod N$ permutes $\{1,\ldots,(N{-}1)/2\}$
when $N$ is odd prime) gives
\[
  \prod_{k=1}^{(N-1)/2}\sin(2\pi k/N)
  = \frac{\sqrt{N}}{2^{(N-1)/2}}.
\]
Therefore
\[
  \prod s_k = \frac{N^{(N-1)/2}}{2^{(N-1)/2}\cdot\sqrt{N}/2^{(N-1)/2}}
  = \frac{N^{(N-1)/2}}{\sqrt{N}} = N^{(N-2)/2}.
\]
For $N{=}11$: $\prod s_k = 11^{9/2} \approx 48558.7$.
Verified numerically for all odd $N$ from $3$ to $37$.
\end{proof}

\subsection{The antisymmetric eigenvalues and a key identity}

\begin{proposition}[Antisymmetric eigenvalues]
\label{prop:a-eigenvalues}
The nonzero eigenvalues of the antisymmetric part~$A$ are purely imaginary:
\begin{equation}\label{eq:a-eigenvalues}
  \Lambda_k(A) = \pm i\,N\,s_k, \qquad k = 1,\ldots,\tfrac{N{-}1}{2},
\end{equation}
where $s_k$ are the eigenvalues of~$S$ from
Theorem~\ref{thm:eigenvalue-closedform}.  In particular,
$|\Lambda_k(A)| = N\,|s_k|$.
\end{proposition}

\begin{proof}
By \eqref{eq:a-circulant}, the matrix $A$ is an antisymmetric circulant
with DFT eigenvalues
$\hat{a}(r) = N\sum_d \hat{f}(d)\,\omega^{rd}$.
Since $\hat{f}$ is odd, these are purely imaginary:
$\hat{a}(r) = -iN\csc(2\pi r/N)/2$.
This gives $\Lambda_k(A) = -iN\,s_k$ for the paired modes.
\end{proof}

\begin{theorem}[Nilpotent identity]
\label{thm:nilpotent-identity}
For the standard order-$N$ Qameah,
\begin{equation}\label{eq:nilpotent-id}
  N^2 S^2 + A^2 = 0.
\end{equation}
\end{theorem}

\begin{proof}
\textbf{Algebraic proof.}
Both $S$ and $A$ are simultaneously diagonalized by the DFT (with
the reversal pairing).  At mode~$k$:
\[
  N^2 s_k^2 + \Lambda_k(A)^2 = N^2 s_k^2 + (-iNs_k)^2
  = N^2 s_k^2 - N^2 s_k^2 = 0.
\]
Since this holds mode-by-mode, $N^2 S^2 + A^2 = 0$ as a matrix identity.

\end{proof}

\begin{corollary}[Antisymmetric trace]
\label{cor:a-trace}
$\mathrm{tr}(A^2) = -N^2\,\mathrm{tr}(S^2) = -N^4(N^2{-}1)/12$.
For $N{=}11$: $\mathrm{tr}(A^2) = -146410$.
\end{corollary}

%=====================================================================
\section{Schr\"odinger Dynamics}
\label{sec:schrodinger}

\begin{theorem}[Schr\"odinger dynamics]
\label{thm:dynamics}
Let $\psi(t) = e^{-iSt}\,\psi_0$ be the Schr\"odinger evolution
generated by~$S$.  Then the observable
\begin{equation}\label{eq:dynamics-observable}
  \mathcal{O}(t) = \langle\psi(t)|A|\psi(t)\rangle
\end{equation}
oscillates at exactly the $(N{-}1)/2$ eigenfrequency pairs
$f_k = |s_k|/\pi$, where $\{s_k\}$ are the nonzero eigenvalues of~$S$.
\end{theorem}

\begin{proof}
Expand $\psi_0$ in the eigenbasis of~$S$: $\psi_0 = \sum_j c_j\,|j\rangle$
where $S|j\rangle = s_j|j\rangle$.  Then
\[
  \mathcal{O}(t) = \sum_{j,k} \overline{c_j}\,c_k\,
  e^{i(s_j - s_k)t}\,\langle j|A|k\rangle.
\]
The anticommutation $\{S,A\}=0$ implies that
$\langle j|A|k\rangle \neq 0$ only when $s_j + s_k = 0$, i.e.\
$A$ connects eigenstates with eigenvalues $+s_k$ and $-s_k$.  For
such a pair, the time-dependent phase is $e^{2is_k t}$, and the
observable oscillates at frequency $f_k = |s_k|/\pi$.

Since $S$ has $(N{-}1)/2$ distinct positive eigenvalues paired with
their negatives, the observable~$\mathcal{O}(t)$ contains exactly
$(N{-}1)/2$ frequency components.  No frequency mixing occurs because
$A$ is strictly off-diagonal in the paired eigenspaces.
\end{proof}

\begin{table}[ht]
\centering
\caption{Eigenfrequency pairs of the $N{=}11$ Qameah symmetric part~$S$.
Frequencies $f_k = s_k/\pi$ in natural units; eigenvalues from
Theorem~\ref{thm:eigenvalue-closedform}.}
\label{tab:eigenfreq}
\begin{tabular}{cccccc}
\toprule
$k$ & 1 & 2 & 3 & 4 & 5 \\
\midrule
$s_k$ & $\frac{11}{2\sin(2\pi/11)}$ {\scriptsize($\approx 19.52$)} & $\frac{11}{2\sin(4\pi/11)}$ {\scriptsize($\approx 10.17$)} & $\frac{11}{2\sin(6\pi/11)}$ {\scriptsize($\approx 7.28$)} & $\frac{11}{2\sin(8\pi/11)}$ {\scriptsize($\approx 6.05$)} & $\frac{11}{2\sin(10\pi/11)}$ {\scriptsize($\approx 5.56$)} \\
$f_k$ & $\frac{11}{2\pi\sin(2\pi/11)}$ {\scriptsize($\approx 6.21$)} & $\frac{11}{2\pi\sin(4\pi/11)}$ {\scriptsize($\approx 3.24$)} & $\frac{11}{2\pi\sin(6\pi/11)}$ {\scriptsize($\approx 2.32$)} & $\frac{11}{2\pi\sin(8\pi/11)}$ {\scriptsize($\approx 1.93$)} & $\frac{11}{2\pi\sin(10\pi/11)}$ {\scriptsize($\approx 1.77$)} \\
\bottomrule
\end{tabular}
\end{table}

\noindent
All five eigenfrequencies are recovered to 14-digit precision for every
tested initial state, confirming the completeness of the spectral
decomposition.

\begin{remark}[Discrete Proca equation]
\label{rem:proca}
The discrete curvature $F_{ijk}=A_{ij}+A_{jk}+A_{ki}$ is quantized,
$F_{ijk}\in\{0,\pm N^2\}$, and the codifferential satisfies
$d^*F=N\,A$, yielding a discrete Proca equation with mass parameter
$m^2=N$ (cf.~the discrete gauge theory of~\cite{kitaev}).
Full proofs appear in the companion paper~\cite{paper3}.
\end{remark}

\subsection{Numerical verification across primes}

\begin{table}[ht]
\centering
\caption{Verification of main identities across odd primes.  All
Frobenius norms $\|{\cdot}\|_F$ are below $10^{-10}$, confirming
exact identities to machine precision.  The product
$\prod s_k = N^{(N-2)/2}$ is verified to 14 significant digits.}
\label{tab:verification}
\begin{tabular}{ccccc}
\toprule
$N$ & $\|SA{+}AS\|_F$ & $\|N^2S^2{+}A^2\|_F$ & $\|L_+^2\|_F$ &
$\mathrm{rank}(L_+)$ \\
\midrule
3  & $<10^{-15}$ & $<10^{-15}$ & $<10^{-15}$ & 1 \\
5  & $<10^{-14}$ & $<10^{-14}$ & $<10^{-14}$ & 2 \\
7  & $<10^{-13}$ & $<10^{-13}$ & $<10^{-13}$ & 3 \\
11 & $<10^{-12}$ & $<10^{-12}$ & $<10^{-12}$ & 5 \\
13 & $<10^{-12}$ & $<10^{-11}$ & $<10^{-11}$ & 6 \\
17 & $<10^{-11}$ & $<10^{-10}$ & $<10^{-10}$ & 8 \\
37 & $<10^{-9}$  & $<10^{-8}$  & $<10^{-8}$  & 18 \\
\bottomrule
\end{tabular}
\end{table}

%=====================================================================
\section{Nilpotent Light-Cone Structure}
\label{sec:nilpotent}

\subsection{Nilpotency and rank}

\begin{definition}[Light-cone operators]
\label{def:lightcone}
For the deviated Qameah with parts $S$ and $A$, define
\begin{equation}\label{eq:lpm}
  L_{\pm} = S \pm \frac{1}{N}\,A.
\end{equation}
\end{definition}

\begin{theorem}[Nilpotency]
\label{thm:nilpotent}
$L_{\pm}^2 = 0$.
\end{theorem}

\begin{proof}
Expanding directly:
\begin{align}
  L_{\pm}^2
  &= \Bigl(S \pm \frac{1}{N}A\Bigr)^2
   = S^2 \pm \frac{1}{N}(SA + AS) + \frac{1}{N^2}A^2 \notag \\
  &= S^2 \pm \frac{1}{N}\{S,A\} + \frac{1}{N^2}A^2.
  \label{eq:lpm-expand}
\end{align}
By Theorem~\ref{thm:anticommute}, $\{S,A\} = 0$, so the middle term
vanishes.  By Theorem~\ref{thm:nilpotent-identity},
$A^2 = -N^2 S^2$, so
\[
  L_{\pm}^2 = S^2 + \frac{1}{N^2}(-N^2 S^2) = S^2 - S^2 = 0.
\]
\end{proof}

\begin{theorem}[Maximal rank]
\label{thm:maxrank}
$\mathrm{rank}(L_+) = (N{-}1)/2$.  This is the maximum possible rank
for a nilpotent-2 operator on $\mathbb{R}^N$ when $N$ is odd.
\end{theorem}

\begin{proof}
For a nilpotent-2 operator $L$ on $\mathbb{R}^N$,
$\operatorname{Im}(L) \subseteq \ker(L)$, so
$\mathrm{rank}(L) \leq \dim\ker(L) = N - \mathrm{rank}(L)$, giving
$\mathrm{rank}(L) \leq \lfloor N/2 \rfloor = (N{-}1)/2$ for odd~$N$.

To show that $L_+$ achieves this bound, we work in the DFT basis.
The operator $L_+$ decomposes into $(N{-}1)/2$ paired $2\times 2$
blocks $B_k$ (modes $k$ and $N{-}k$) plus one zero mode ($k=0$).
At mode~$k$:
\[
  B_k = \begin{pmatrix} s_k & a_k/N \\ -a_k/N & -s_k \end{pmatrix}
\]
where $s_k = N/(2\sin(2\pi k/N))$ and $a_k = iNs_k$
(Proposition~\ref{prop:a-eigenvalues}).  Then $|a_k/N| = |s_k|$, so
$\det(B_k) = -s_k^2 + s_k^2 = 0$ and $\mathrm{rank}(B_k) = 1$.
Summing: $\mathrm{rank}(L_+) = (N{-}1)/2$.

Verified numerically for all odd primes $N = 3, 5, 7, 11, 13, 17, 19,
23, 29, 31, 37$.
\end{proof}

\begin{corollary}[Lagrangian subspace]
\label{cor:lagrangian}
$\operatorname{Im}(L_+) \subseteq \ker(L_+)$ with
$\dim\operatorname{Im}(L_+) = (N{-}1)/2$
and $\dim\ker(L_+) = (N{+}1)/2$.
The image is a maximal totally isotropic subspace of $\mathbb{R}^N$
under the nilpotent bilinear form $B(x,y) = x^T L_+ y$: a natural
\emph{Lagrangian subspace}.
\end{corollary}

\subsection{BRST cohomology}

The nilpotent operator $L_+$ defines a cochain complex
$0 \to \mathbb{R}^N \xrightarrow{L_+} \mathbb{R}^N \to 0$, and
the cohomology $H^1(L_+) = \ker(L_+)/\operatorname{Im}(L_+)$ measures
the space of ``gauge-invariant'' observables~\cite{becchi,henneaux}.

\begin{theorem}[One-dimensional cohomology]
\label{thm:brst}
For all standard Qameah of odd order~$N$,
\begin{equation}\label{eq:brst-dim}
  \dim H^1(L_+) = \dim\frac{\ker(L_+)}{\operatorname{Im}(L_+)} = 1.
\end{equation}
The unique cohomology class is generated by the all-ones vector
$\mathbf{1} = (1,1,\ldots,1)^T$.
\end{theorem}

\begin{proof}
\textbf{Dimension count.}  From Theorem~\ref{thm:maxrank},
$\dim\ker(L_+) = (N{+}1)/2$ and $\dim\operatorname{Im}(L_+) = (N{-}1)/2$.
Therefore $\dim H^1(L_+) = (N{+}1)/2 - (N{-}1)/2 = 1$.

\medskip\noindent
\textbf{Generator identification.}  We show
$\mathbf{1} \in \ker(L_+) \setminus \operatorname{Im}(L_+)$.
\begin{itemize}
\item $\mathbf{1} \in \ker(L_+)$: the rows of $L_+ = S + A/N$ sum to
zero.  Indeed, $S$ has zero row sums (deviated magic square) and $A$
has zero row sums (antisymmetric matrix), so
$L_+\,\mathbf{1} = \mathbf{0}$.

\item $\mathbf{1} \notin \operatorname{Im}(L_+)$: the columns of $L_+$
also sum to zero ($S$ has zero column sums by the magic property, and
$A$ has zero column sums since $A^T = -A$ with zero row sums).
Therefore $\mathbf{1}^T L_+ = \mathbf{0}^T$, which means
$\mathbf{1}$ is orthogonal to every vector in $\operatorname{Im}(L_+)$.
Since $\mathbf{1} \neq \mathbf{0}$, it cannot lie in
$\operatorname{Im}(L_+)$.
\end{itemize}
Hence $[\mathbf{1}]$ generates $H^1(L_+)$.
Verified for $N = 3, 5, 7, \ldots, 37$ (13 values).
\end{proof}

\begin{remark}[Quantum error-correcting codes]
\label{rem:qec}
The nilpotency $L_+^2 = 0$ (Theorem~\ref{thm:nilpotent})
implies that the row space of $L_+ \bmod N$ is self-orthogonal over $\mathbb{F}_N$.
By the CSS construction, this yields a quantum error-correcting code
$[[N, N-4, 3]]_N$ for every prime $N \geq 5$.
This code saturates the quantum Singleton bound and is therefore
a quantum MDS code.
The generator matrix is a Reed--Solomon $[N,2,N-1]_N$ code,
with the $\mathrm{GL}(2,\mathbb{Z}_N)$ magic square entries as evaluation points.
\end{remark}

%=====================================================================
\section{Three-Generation Equidistribution for \texorpdfstring{$N=11$}{N=11}}
\label{sec:three-gen}

\subsection{Simplicial face counting}

The $N{=}11$ Qameah acts on $\mathbb{R}^{11}$, whose standard
basis vectors $\{e_0, e_1, \ldots, e_{10}\}$ span a 10-simplex
$\Delta^{10}$~\cite{coxeter}.  The triangular (2-dimensional) faces of $\Delta^{10}$
are indexed by 3-element subsets of $\{0,1,\ldots,10\}$:
\begin{equation}\label{eq:face-count}
  \binom{11}{3} = 165.
\end{equation}
For each face $\{i,j,k\}$, define the \emph{$S$-weight} as the
sum of the absolute values of the relevant entries of~$S$:
\begin{equation}\label{eq:s-weight}
  w_S(\{i,j,k\}) = 2(|S_{ij}| + |S_{ik}| + |S_{jk}|).
\end{equation}
The total $S$-weight over all 165 faces is
$W_{\mathrm{tot}} = 3 \times 20130 = 60390$.

\subsection{Balanced partitions}

\begin{definition}[Balanced 3-partition]
A \emph{balanced 3-partition} of the 165 faces is a decomposition
$\mathcal{F} = \mathcal{G}_1 \sqcup \mathcal{G}_2 \sqcup \mathcal{G}_3$
with $|\mathcal{G}_g| = 55$ for each $g \in \{1,2,3\}$ and
\begin{equation}\label{eq:weight-balance}
  \sum_{F \in \mathcal{G}_g} w_S(F) = 20130 \quad
  \text{for each } g.
\end{equation}
\end{definition}

\begin{theorem}[Equidistribution]
\label{thm:three-gen}
There exist exactly $16 = 2^4$ balanced 3-partitions of the 165
triangular faces of $\Delta^{10}$ satisfying~\eqref{eq:weight-balance}.
\end{theorem}

The proof is constructive (exhaustive enumeration).  The 165 faces
decompose into orbits under the action of the cyclic group
$\mathbb{Z}_{11}$ (generated by the index shift $i \mapsto i+1 \bmod 11$).
Each orbit has size~11, giving $165/11 = 15$ orbits.  Of these:

\begin{itemize}
\item \textbf{5 self-conjugate orbits:} These are fixed under the
      complement map $\{i,j,k\} \mapsto \{N{-}i, N{-}j, N{-}k\}$.
      They form Generation~1 (55 faces).

\item \textbf{10 paired orbits:} The remaining orbits come in 5
      conjugate pairs under the complement map.  Assigning one partner
      to Generation~2 and the other to Generation~3 gives a balanced
      partition, provided the $S$-weights match.

\item \textbf{$2^4$ choices:} Of the 5 conjugate pairs, 4 admit a
      free binary choice (either partner can go to either generation
      without breaking the weight balance), and 1 pair is fixed by
      the constraint.  This gives $2^4 = 16$ solutions.
\end{itemize}

\noindent
The enumeration code is available at
\url{https://github.com/Insightful-Calculations/qhots_v7_calc}.

\begin{proposition}[Edge balance]
\label{prop:edge-balance}
In every balanced 3-partition, each generation $\mathcal{G}_g$ contains
exactly 55 edges (1-dimensional faces), and every edge participates
in exactly $c_g(e) = 3$ triangular faces within its generation.
\end{proposition}

\begin{remark}[Random hit rate]
A random 3-partition of 165 objects into groups of 55 achieves
exact weight balance $20130/20130/20130$ with probability
$\approx 0.035\%$.  The existence of 16 such partitions with
additional edge-balance and orbit structure is a nontrivial
combinatorial property of the $N{=}11$ Qameah.
\end{remark}

\begin{remark}[Universality for other $N$]
Whether analogous balanced 3-partitions exist for other odd primes
is an open question.  The structure depends on the specific values
of $\binom{N}{3}$, the orbit decomposition under $\mathbb{Z}_N$,
and the $S$-weight distribution.  The case $N{=}11$ is distinguished
by $\binom{11}{3} = 165 = 3 \times 55$, where $55 = T_{10}$
(the 10th triangular number), enabling the threefold split.
For $N{=}7$: $\binom{7}{3} = 35$, which is not divisible by~3,
so no balanced 3-partition exists.  For $N{=}13$:
$\binom{13}{3} = 286$, which is not divisible by~3 either.
Hence the three-generation structure is specific to $N{=}11$
among the first several primes.
\end{remark}

%=====================================================================
\section{Scope of the Results}
\label{sec:honest}

\subsection{What is proved}

\begin{enumerate}
\item \textbf{Anticommutation $\{S,A\}=0$}
      (Theorem~\ref{thm:anticommute}): Proved algebraically for
      all standard Qameah of odd order, with the proof
      depending only on the complementarity of the circulant generator.

\item \textbf{Eigenvalue closed form}
      (Theorem~\ref{thm:eigenvalue-closedform}): Proved via the
      Hankel--circulant factorization $S = J_{\mathrm{rev}}\,C_h$ and
      standard DFT eigenvalue formulas.  Universal for all odd~$N$.

\item \textbf{Nilpotent identity $N^2S^2 + A^2 = 0$}
      (Theorem~\ref{thm:nilpotent-identity}): Proved mode-by-mode in
      the DFT basis.  Universal for all odd~$N$.

\item \textbf{$L_\pm^2 = 0$ and maximal rank}
      (Theorems~\ref{thm:nilpotent}--\ref{thm:maxrank}): Direct
      consequences of (1)--(3).  The rank bound is tight.

\item \textbf{BRST cohomology $\dim H^1 = 1$}
      (Theorem~\ref{thm:brst}): Follows from the dimension count.
      Universal for all odd~$N$.

\item \textbf{Spectral generating function}
      (Theorem~\ref{thm:t95}): Proved for all odd~$N$; closes the
      spectral arithmetic series T91--T97.

\item \textbf{Three-generation equidistribution}
      (Theorem~\ref{thm:three-gen}): Constructive proof by exhaustive
      enumeration, specific to $N{=}11$.
\end{enumerate}

\subsection{What is not proved}

\begin{enumerate}
\item \textbf{No physical content is claimed.}  The
      ``Schr\"odinger dynamics'' of Theorem~\ref{thm:dynamics} is a
      mathematical statement about the spectral decomposition of
      $e^{-iSt}$, not a claim about quantum mechanics.  The
      ``light-cone operators'' $L_\pm$ are named by analogy; no
      physical light cone is implied.

\item \textbf{No coupling to physics.}  Whether the algebraic
      structures identified here (anticommutation, nilpotency, BRST
      cohomology) have physical realizations in quantum field theory,
      condensed matter, or any other domain is an open question not
      addressed in this paper.

\item \textbf{No E8 connection is derived.}  The coincidence that
      $(N^2{-}1)/4 = 30 = h(E_8)$ (the Coxeter number of the
      exceptional Lie algebra $E_8$) at $N{=}11$ is observed but
      not explained.  It may be accidental.

\item \textbf{Strengthened divisibility.}  The bound
      $N^{4k} \mid e_k$ (stronger than Theorem~\ref{thm:t91}'s
      $N^{2k} \mid e_k$) is verified only for $N{=}11$ and lacks a
      general proof.

\item \textbf{Composite $N$.}  All theorems are stated for odd
      prime~$N$.  Composite odd~$N$ introduces additional
      divisibility structure that may require modifications to the
      generating function.
\end{enumerate}

\subsection{Negative results}

\begin{enumerate}
\item Attempts to derive physical coupling constants from the
      eigenvalue ratios $s_j/s_k$ found no matches beyond the level
      of numerical coincidence.

\item The Weinberg angle $\sin^2\theta_W$ cannot be obtained from
      any ratio of $S$-eigenvalues at $N{=}11$.

\item A discrete Proca Lagrangian $\mathcal{L}=N^5(N^2{-}1)/96$ with $d^*F=N\,A$
      has been identified (Remark~\ref{rem:proca}); the $\{S,A\}=0$
      relation itself is not yet derived from this action.

\item The three-generation structure does not extend to $N{=}7$
      or $N{=}13$ (the adjacent primes), due to the divisibility
      constraint $3 \mid \binom{N}{3}$.

\item Although the anticommutation $\{S,A\}=0$ is proved algebraically
      for all odd~$N$ (Theorem~\ref{thm:anticommute}), analogous
      closed-form proofs of the strengthened spectral divisibility
      $N^{4k}\mid e_k$ have not been found.
\end{enumerate}

%=====================================================================
\section{Open Problems}
\label{sec:limitations}

\begin{enumerate}
\item \textbf{Strengthened divisibility for general $N$.}
      Theorem~\ref{thm:t91} establishes $N^{2k} \mid e_k$ for all odd
      primes.  The stronger bound $N^{4k} \mid e_k$, verified for $N=11$,
      lacks a general proof.  A $p$-adic or motivic approach may be
      required.

\item \textbf{Classification of nilpotent magic squares.}
      Non-Qameah members of the de~la~Loub\`ere family violate $L_+^2 = 0$ for $N\ge 5$
      (Section~\ref{sec:nilpotent}).  A complete classification of which
      magic-square constructions yield the nilpotent light-cone structure
      is open.  The relevant condition appears to be the complementarity
      of the circulant generator (Lemma~\ref{lem:complementarity}).

\item \textbf{Composite odd $N$.}
      All theorems are stated for odd prime~$N$.  For composite $N$
      (e.g.\ $N = 9, 15, 21$), the DFT block structure changes and the
      generating function (Theorem~\ref{thm:t95}) may require
      modification.  Preliminary numerics suggest that anticommutation
      and nilpotency persist, but the rank formula changes.

\item \textbf{Three-generation structure beyond $N=11$.}
      The balanced 3-partition of Theorem~\ref{thm:three-gen} is specific
      to $N=11$.  The obstruction for other primes is the divisibility
      $3 \mid \binom{N}{3}$, which fails for $N = 7, 13, 17$.
      Whether weaker forms of equidistribution exist is unexplored.

\item \textbf{From discrete Proca to $\{S,A\}=0$.}
      A discrete Proca Lagrangian $\mathcal{L}=N^5(N^2{-}1)/96$ now exists
      (Remark~\ref{rem:proca}), but the anticommutation $\{S,A\}=0$
      is not yet derived as its equation of motion.  Whether the discrete Proca
      action \emph{implies} $\{S,A\}=0$ or requires additional structure
      remains open.

\item \textbf{Connections to representation theory.}
      The coincidence that each row of $|S|$ sums to $(N^2{-}1)/4$, which
      equals the Coxeter number $h(E_8) = 30$ at $N{=}11$, suggests a
      possible link to exceptional Lie algebras.  No such connection has
      been established.
\end{enumerate}

%=====================================================================
\begin{acknowledgments}
This paper is the product of a human--AI collaboration.
Hr\'oar \TH\'or Reynisson conceived the Qameah research programme,
directed all investigations, and validated the mathematical content.
Claude (Anthropic) contributed proof construction, algebraic verification,
and manuscript drafting as a co-author.
Computational verification was performed using NumPy/SciPy
and the Sefirot autonomous research system.
All numerical results reported have been verified for odd primes
$N = 3, 5, 7, 11, 13, 17, 19, 23, 29, 31, 37$.
\end{acknowledgments}

%=====================================================================
\begin{thebibliography}{99}

\bibitem{coxeter} H.S.M. Coxeter, \textit{Regular Polytopes}, 3rd ed.
(Dover, New York, 1973).

\bibitem{davis} P.J. Davis, \textit{Circulant Matrices}, 2nd ed.
(Chelsea, New York, 1994).

\bibitem{andrews} W.S. Andrews, \textit{Magic Squares and Cubes},
2nd ed. (Dover, New York, 1960).

\bibitem{becchi} C. Becchi, A. Rouet, and R. Stora, ``Renormalization of
gauge theories,'' Ann. Phys. \textbf{98}, 287 (1976).

\bibitem{chebyshev} P.L. Chebyshev, ``Sur les expressions approximatives
des int\'egrales d\'efinies par les autres prises entre les m\^emes
limites,'' Proc. Acad. Sci. St. Petersburg (1882).

\bibitem{deLaLoubere} S. de la Loub\`ere, \textit{Du Royaume de Siam}
(J.-B. Coignard, Paris, 1691), Vol.~2.

\bibitem{kraitchik} M. Kraitchik, \textit{Mathematical Recreations},
2nd ed. (Dover, New York, 1953).

\bibitem{terras} A. Terras, \textit{Fourier Analysis on Finite Groups
and Applications} (Cambridge Univ. Press, 1999).

\bibitem{henneaux} M. Henneaux and C. Teitelboim, \textit{Quantization
of Gauge Systems} (Princeton Univ. Press, 1992).

\bibitem{gray} R.M. Gray, ``Toeplitz and Circulant Matrices: A Review,''
Found. Trends Commun. Inf. Theory \textbf{2}, 155 (2006).

\bibitem{appleby} D.M. Appleby, ``Symmetric informationally complete
positive operator valued measures and the extended Clifford group,''
J. Math. Phys. \textbf{46}, 052107 (2005).

\bibitem{nordgren} R.P. Nordgren, ``On properties of special magic square
matrices,'' Linear Algebra Appl. \textbf{437}, 2009 (2012).

\bibitem{washington} L.C. Washington, \textit{Introduction to Cyclotomic
Fields}, 2nd ed. (Springer, New York, 1997).

\bibitem{kitaev} A.Yu. Kitaev, ``Fault-tolerant quantum computation by
anyons,'' Ann. Phys. \textbf{303}, 2 (2003).

\bibitem{mattingly} R.B. Mattingly, ``Even order regular magic squares are
singular,'' Amer. Math. Monthly \textbf{107}, 777 (2000).

\bibitem{paper3} H. \TH. Reynisson and Claude Opus 4.6,
``Clifford algebraic and gauge-theoretic structure of the Qameah,''
in preparation (2026).

\end{thebibliography}

\end{document}
